% Faithful transcription — original English. Transcribed from the original first printing:
% Arthur Cayley, "On the Deficiency of Certain Surfaces", Mathematische Annalen, Band 3,
% Heft 4 (1871), pp. 526–529. Scan: EuDML doc/156505 (printed pp. 526–529 = leaves n534–n537).
%
% \origpage uses the PRINTED page numbers (526–529). Running heads and page numbers are omitted.
%
% Cross-page typography and rendering decisions:
% - Cuspidal lines on pp. 526–527 use Greek kappa \varkappa, matching the 19th-century Teubner
%   mathematical font, distinguished from Latin x. On p. 528, x denotes Salmon's space curve double points.
% - Coordinates on p. 529 use Latin x, y, z and Greek \omega.
% - German quotation marks „..." are preserved as literal Unicode per HOUSESTYLE R22.
% - Printer's errors faithfully reproduced per HOUSESTYLE R4:
%     p. 527: "wich" without "h" ("the nature of wich I do not completely understand").
%     p. 528: double hyphen "--" in "10m -- 3n".
%     p. 529: double hyphen "--" in "\delta y -- \beta\omega".
% See notation.md for the complete specification.
\documentclass{article}
\usepackage{readmasters}
\begin{document}

\origpage{526}
\section*{On the Deficiency of Certain Surfaces.}

\subsection*{By A.~Cayley.}

If a given point or curve is to be an ordinary or singular point or curve on a surface of the order $n$, this imposes on the surface a certain number of conditions, which number may be termed the „Postulation“; thus „Postulation of a given curve \emph{quà} $i$-tuple curve on a surface $n$“ will denote the number of conditions to be satisfied by the surface in order that the given curve may be an $i$-tuple curve on the surface.

The „deficiency“ (Flächengeschlecht) of a given surface of the order $n$ is
\[
= \frac{1}{6} (n - 1) (n - 2) (n - 3) \text{ less deficiency-value of the several singularities}
\]
viz.\ as shown by Dr.~Noether, if the surface has a given $i$-tuple curve, the deficiency-value hereof is
\[
= \text{Postulation of the curve \emph{quà} } (i - 1) \text{ tuple curve on a surface } n - 4\text{;}
\]
and if the surface has an $i$-conical point, the deficiency-value hereof is
\[
= \text{Postulation of the point \emph{quà} } (i - 2) \text{ conical point (on a surface } n - 4\text{); viz.\ this is } = \frac{1}{6} i(i - 1)(i - 2)\text{,}
\]
and is thus independent of the order of the surface.

I remark that if the tangent-cone at the $i$-conical point has $\delta$ double lines and $\varkappa$ cuspidal lines, then the deficiency-value is
\[
= \frac{1}{6} i(i - 1)(i - 2) + (i - 2)(n - i - 1)(\delta + \varkappa).
\]
In the case of a double or cuspidal curve $i$ is $= 2$, and the deficiency-value is
\[
= \text{Postulation of given curve \emph{quà} simple curve on a surface } n - 4\text{;}
\]
and so for an ordinary conical point $i$ is $= 2$, and the deficiency-value is $= 0$: results which were first obtained by Dr.~Clebsch.

I found in this manner the expression for the deficiency of a surface $n$ having a double and cuspidal curve and the other singularities considered in my „Memoir on the theory of Reciprocal surfaces“ Phil.\ Trans.\ vol.\ 159 (1869); viz.\ this was
\origpage{527}
\[
D = \frac{1}{6}(n - 1)(n - 2)(n - 3) - (n - 3)(b + c) + \frac{1}{2}(q + r) + 2t + \frac{1}{2}\beta + \frac{1}{2}\gamma + i - \frac{1}{2}\Theta,
\]
where we have\ednote{In the explanation of $\Theta$ in the list below, the printed text reads ``the nature of wich I do not completely understand'': so printed, without an ``h'', for ``which''.}
\[
\begin{array}{ll}
b, & \text{order of double curve,} \\
q, & \text{class of Do.,} \\
c, & \text{order of cuspidal curve,} \\
r, & \text{class of Do.,} \\
\beta, & \text{number of intersections of the two curves, stationary points} \\
& \quad\text{on } b, \\
\gamma, & \text{number of intersections, stationary points on } c, \\
i, & \text{number of intersections, not stationary points on either curve,} \\
\Theta, & \text{number of certain singular points on } c\text{, the nature of wich} \\
& \quad\text{I do not completely understand, and which is here taken to} \\
& \quad\text{be } = 0.
\end{array}
\]

Before going further I remark that
\[
\begin{aligned}
\text{Postulation of right line \emph{quà} } i\text{-tuple on surface } n & \\
&= \frac{1}{2} i(i + 1)n - \frac{1}{6} i(i + 1)(2i - 5), \\
&= \frac{1}{6} i(i + 1)(3n - 2i + 5).
\end{aligned}
\]
Whence if a surface $n$ has a $i$-tuple right line, the deficiency-value hereof is
\[
= \frac{1}{6} i(i - 1)(3n - 2i - 5),
\]
or we have
\[
\begin{aligned}
D &= \frac{1}{6}(n - 1)(n - 2)(n - 3) - \frac{1}{6} i(i - 1)(3n - 2i - 5) \\
&= \frac{1}{6}(i - n + 1)(i - n + 2)(2i + n - 3);
\end{aligned}
\]
so that $D = 0$ if either $i = n - 1$ or $i = n - 2$; the former case is that of a scroll (skew surface) with a $(n - 1)$ tuple right line the latter that of a surface with a $(n - 2)$ tuple line: whence (as shown by Dr.~Noether) such surface is rationally transformable into a plane.

For a surface of the order $n$ with a $i$-conical point where the tangent cone has $\delta$ double lines and $\varkappa$ cuspidal lines, we have
\[
\begin{aligned}
D &= \frac{1}{6}(n - 1)(n - 2)(n - 3) - \{\frac{1}{6} i(i - 1)(i - 2) + (i - 2)(n - i - 1)(\delta + \varkappa)\} \\
&= \frac{1}{6}(n - i - 1)\{n^2 + n(i - 5) + i^2 - 4i + 6 - 6(i - 2)(\delta + \varkappa)\}
\end{aligned}
\]
viz.\ for $i = n - 1$ this is $D = 0$ (in fact a surface $n$ with a $(n - 1)$ conical point is at once seen to be rationally transformable into a plane): and for $i = n$, that is for a cone of the order $n$, we have
\[
D = -\frac{1}{2}(n - 1)(n - 2) + (n - 2)(\delta + \varkappa) - (n - 3)(\delta + \varkappa),
\]
where the last term $-(n - 3)(\delta + \varkappa)$ is added because in the present case the surface has the $\delta$ double lines and the $\varkappa$ cuspidal lines.

The formula therefore gives
\[
D = -\frac{1}{2}(n - 1)(n - 2) + \delta + \varkappa,
\]
viz.\ this is equal to the deficiency of the plane sections \emph{taken negatively}.

\origpage{528}
I find that the same property exists \emph{first} in the case of a scroll (skew surface) having only a double curve; and \emph{secondly} in the case of a torse (developable surface) having a cuspidal curve with the ordinary singularities; and this being so there can I think be no doubt but that it is true for any scroll or torse whatever --- viz.\ that for any ruled surface whatever the deficiency is equal to that of the plane section taken negatively.

First for the scroll we have
\[
D = \frac{1}{6}(n - 1)(n - 2)(n - 3) - (n - 3)b + \frac{1}{2}q + 2t,
\]
which should be
\[
= -\frac{1}{2}(n - 1)(n - 2) + b.
\]
Salmon's equations give in the case of a scroll
\[
\begin{aligned}
3t &= (n - 4)\{3b - n(n - 2)\}, \\
q &= n(n - 2)(n - 5) - 2(n - 6)b,
\end{aligned}
\]
and with these values the relation is at once verified.

Secondly for the torse; changing the notation into that used for the singularities of the curve and torse, we have
\[
D = \frac{1}{6}(r - 1)(r - 2)(r - 3) - (r - 3)(x + m) + \frac{1}{2}(q + r) + 2t + \frac{7}{2}\beta + \frac{5}{2}\gamma + \alpha,
\]
which should be
\[
= -\frac{1}{2}(m - 1)(m - 2) + h + \beta;
\]
we have $q = r(n - 3) - 3\alpha$, and substituting this value and expressing every thing in terms of $r, m, n$ by means of the formula\ednote{In the expression for $h$ in the following system, the print sets a double hyphen between $10m$ and $3n$, $10m -- 3n$: reproduced here as printed.}
\[
\begin{aligned}
x &= \frac{1}{2}(r^2 - r - n - 3m), \\
\alpha &= m - 3r + 3n, \\
\beta &= n - 3r + 3m, \\
t &= \frac{1}{6}\{r^3 - 3r^2 - 58r - 3r(n + 3m) + 42n + 78m\}, \\
\gamma &= rm + 12r - 14m - 6n, \\
h &= \frac{1}{2}(m^2 - 10m -- 3n + 8r),
\end{aligned}
\]
we have after all reductions\ednote{In the rightmost member of the following display, the minus sign in $(m - 1)$ shows a small ink defect above the rule in the print, resembling a division sign at first glance; it is a defective minus sign, matching the identical factor $(m - 1)$ earlier on the same page.}
\[
D = -\frac{1}{2}(m + n) + r - 1 = -\frac{1}{2}(m - 1)(m - 2) + h + \beta.
\]

We have thus a class of surfaces of \emph{negative deficiency}; viz.\ any rational transformation of a cone for which the plane section has a given (positive) deficiency produces such a surface: and I think it may be assumed conversely that a surface of negative deficiency is always the rational transformation of a cone for which the deficiency is equal to that of the surface taken with the reverse sign. As an instance take a quintic surface having a nodal conic and two 3-conical (cubiconical) points (this of course implies that the line joining the two cubiconical points is a line on the surface): the formula for the
\origpage{529}
\[
D = \frac{1}{6}(n - 1)(n - 2)(n - 3) - (n - 3)b + \frac{1}{2}q + 2t - 2
\]
(viz.\ a term $-1$ for each of the cubiconical points)
\[
= 4 - 4 + 1 - 2, \quad = -1.
\]

Such a surface can be obtained as the quadric inverse of a cubic cone; viz.\ taking for the vertex the point $x : y : z : \omega = \alpha : \beta : \gamma : \delta$ and the cone to pass thro' the point $x = 0, y = 0, z = 0$, the equation of the cone is\ednote{In the following equation the print sets a double hyphen between $\delta y$ and $\beta\omega$, $\delta y -- \beta\omega$, matching the typographical slip in the formula for $h$ on the preceding page: reproduced here as printed.}
\[
(\delta x - \alpha\omega, \, \delta y -- \beta\omega, \, \delta z - \gamma\omega)^3 = 0
\]
where $(\alpha, \beta, \gamma)^3 = 0$.

Taking $Q$ a quadric function $(x, y, z)^2$, the transformation in question consists in the change of $x, y, z, \omega$ into $x\omega, y\omega, z\omega, Q$; viz.\ the new equation, rejecting the factor $\omega$ which divides out, is
\[
\frac{1}{\omega} (\delta x\omega - \alpha Q, \, \delta y\omega - \beta Q, \, \delta z\omega - \gamma Q)^3 = 0
\]
which is a quintic surface having the two cubiconical points $x = 0, y = 0, z = 0$ and $x : y : z : \omega = \alpha : \beta : \gamma : \frac{1}{\delta} Q_0$ (where $Q_0$ is the value of $Q$ on writing therein $\alpha, \beta, \gamma$ in place of $x, y, z$): and having the nodal conic $\omega = 0, Q = 0$.

Cambridge, 5.\ Jan.\ 1871.

\end{document}
